\section{A critical finite-difference hierarchy of every harmonic rank}
\label{hier:sec}

The original finite-difference construction uses a stencil of order
$2r-2$ and confines its special-function dependence to one moment entry.
Shorter stencils give an explicit hierarchy whose harmonic ranks run
through every value from one to $r$, while the integer-order weights
retain the same critical decay.  Its spectral continuation has a
completely determined staircase of pole orders.  This is a spectral
finite-difference family; the more general overlapping resonances of
the earlier Gauss ensemble are a separate question.

Fix integers
\begin{equation}\label{hier:parameters}
 r\ge2,\qquad 1\le\rho\le r,\qquad
 d=2r-2,\qquad h=d-\rho+1=2r-\rho-1.
\end{equation}
For $\Re a>0$, use the right-half-plane logarithm to define
\begin{equation}\label{hier:weight}
 w_n^{[h]}(s,a)=\frac1{h!}\sum_{j=0}^h(-1)^j\binom hj
                                      (n+a+j)^{-s}.
\end{equation}
The corresponding ordered ensemble is
\begin{equation}\label{hier:ensemble}
 Z_{r,\rho}(s,a)=
 \sum_{0\le n_1<\cdots<n_r}
 \Delta(n_1,\ldots,n_r)^2\prod_{i=1}^r w_{n_i}^{[h]}(s,a).
\end{equation}
There is no additional $1/r!$ because the indices are strictly ordered.

\subsection{Finite moments and exact shifted-zeta closure}

Define finite polynomials and endpoint functions by
\begin{align}
 b_k(a;h)&=\frac1{h!}\sum_{j=0}^h(-1)^j\binom hj(-a-j)^k,
                                                        \label{hier:b}\\
 E_k(s,a;h)&=-\frac1{h!}\sum_{j=1}^h(-1)^j\binom hj
       \sum_{\ell=0}^{j-1}(\ell-j)^k(\ell+a)^{-s}.
                                                        \label{hier:E}
\end{align}
The convention $0^0=1$ applies to zeroth moments.  Polynomial
annihilation gives
\begin{equation}\label{hier:annihilation}
 b_0=\cdots=b_{h-1}=0,\qquad b_h=1.
\end{equation}
For $0\le\ell\le\rho-1$ let $C_\ell(a;h)$ be the symmetric Hankel
matrix with entries
\begin{equation}\label{hier:C}
 (C_\ell)_{ij}=\binom{i+j}{\ell}b_{i+j-\ell}(a;h),
 \qquad 0\le i,j<r,
\end{equation}
with a term set to zero if its lower index is negative or if
$\ell>i+j$.  Also let $E(s,a;h)=(E_{i+j}(s,a;h))_{i,j<r}$.

\begin{theorem}[A finite shifted-zeta determinant]
\label{hier:closure}
The ensemble \eqref{hier:ensemble} is absolutely and locally normally
convergent for $\Re s>\rho$, $\Re a>0$, and there
\begin{equation}\label{hier:closed}
 \boxed{\displaystyle
 Z_{r,\rho}(s,a)=\det\left(
 E(s,a;h)+\sum_{\ell=0}^{\rho-1}
                    \zeta(s-\ell,a)C_\ell(a;h)\right).}
\end{equation}
The right side gives its meromorphic continuation to all $s\in\mathbb C$.
Its only possible poles are $1,2,\ldots,\rho$.
\end{theorem}
\begin{proof}
For $\Re s>0$ the weight has the integral representation
\begin{equation}\label{hier:laplace}
 w_n^{[h]}(s,a)=\frac1{h!\Gamma(s)}
 \int_0^\infty t^{s-1}e^{-(n+a)t}(1-e^{-t})^h\,dt.
\end{equation}
Repeated finite differences, or expansion of this kernel at zero,
give locally uniformly
\begin{equation}\label{hier:decay}
 w_n^{[h]}(s,a)=\frac{(s)_h}{h!}(n+a)^{-s-h}
                     +O((n+1)^{-\Re s-h-1}).
\end{equation}
For ordered indices, $|\Delta(\boldsymbol n)|^2\le
\prod_i(n_i+1)^{2(i-1)}$.  The least decaying factor is summable
when $\Re s+h-d>1$, precisely $\Re s>\rho$.

For a moment of degree $k\le d$, initially take $\Re s>d+1$ so
that all separately shifted sums converge.  Shift the index by $j$
and expand
\[
 (n-j)^k=\sum_{\ell=0}^k\binom k\ell
                      (-a-j)^{k-\ell}(n+a)^\ell.
\]
The lower endpoints give \eqref{hier:E}, and the infinite shifted
sums give
\begin{equation}\label{hier:moment}
 \sum_{n\ge0}n^k w_n^{[h]}(s,a)
 =E_k(s,a;h)+\sum_{\ell=0}^{k-h}
     \binom k\ell b_{k-\ell}(a;h)\zeta(s-\ell,a).
\end{equation}
An empty sum is zero.  Both sides are holomorphic on their common
normal-convergence half-plane, so the identity theorem extends this
identity to $\Re s>\rho$.  Finite Cauchy--Binet and passage to the
normal limit give \eqref{hier:closed}.  Each endpoint function is
entire in $s$, and every Hurwitz factor has just its displayed pole.
\end{proof}

For later cutoff use, define
$H_N^{(z)}(a)=\sum_{n=0}^N(n+a)^{-z}$ for every complex $z$ and
\begin{equation}\label{hier:upper}
 T_{k,N}(s,a;h)=\frac1{h!}\sum_{j=1}^h(-1)^j\binom hj
       \sum_{\ell=N+1}^{N+j}(\ell-j)^k(\ell+a)^{-s}.
\end{equation}
The exact finite counterpart of \eqref{hier:moment} is
\begin{equation}\label{hier:finite-moment}
 \sum_{n=0}^N n^k w_n^{[h]}(s,a)
 = E_k+T_{k,N}+
   \sum_{\ell=0}^{k-h}\binom k\ell b_{k-\ell}
                                     H_N^{(s-\ell)}(a).
\end{equation}
It holds for every $s$, including orders at which the infinite
ensemble has a pole.  It follows from the finite index-shift identity;
no divergent series is split in its proof.

\subsection{Exact rank and the unique highest harmonic monomial}

\begin{theorem}[Every rank at the same critical decay]
\label{hier:rank}
At $s=\rho$ and real $a>0$, the rational weight
$R_{r,\rho,a}(n)=w_n^{[h]}(\rho,a)$ is positive and obeys
\begin{equation}\label{hier:critical}
 R_{r,\rho,a}(n)=\binom{2r-2}{\rho-1}n^{-(2r-1)}
                                      +O(n^{-2r}).
\end{equation}
Its harmonic pencil has generic rank exactly $\rho$.

Put $L=r-\rho$ and let
\begin{equation}\label{hier:Vcutoff}
 V_{L,N}(a)=\det\left(
   \sum_{n=0}^N n^{i+j}R_{r,\rho,a}(n)\right)_{0\le i,j<L},
 \qquad V_{0,N}=1.
\end{equation}
This is rational in $N$ (and rational in $a$).  The degree-$\rho$
homogeneous part of the full $r$-particle cutoff determinant, with
formal harmonic variables $X_q$ corresponding to $H_N^{(q)}(a)$, is
exactly
\begin{equation}\label{hier:top}
 \boxed{\displaystyle
 (-1)^{\rho(\rho-1)/2}V_{L,N}(a)X_\rho^\rho.}
\end{equation}
It is nonzero whenever $N\ge L-1$, with this cutoff condition vacuous
when $L=0$.
\end{theorem}
\begin{proof}
Positivity follows from \eqref{hier:laplace}.  In
\eqref{hier:decay}, $h+\rho=2r-1$ and
$(\rho)_h/h!=\binom{2r-2}{\rho-1}$.

By \eqref{hier:annihilation}, $C_\ell$ is supported in its final
$\rho-\ell$ rows and columns.  In its final
$(\rho-\ell)\times(\rho-\ell)$ block, every entry below the first
possible anti-diagonal in moment order vanishes, and the entries on
that anti-diagonal all equal
$\binom{h+\ell}{\ell}$.  In particular, all matrices lie in the
same final $\rho\times\rho$ block.  Within that block $C_0$ has
zeros for relative indices $i+j<\rho-1$ and ones when
$i+j=\rho-1$; all other $C_\ell$ start strictly after that
anti-diagonal.  Hence
\begin{equation}\label{hier:block-det}
 \det\left(\sum_{\ell=0}^{\rho-1}Y_\ell
                  (C_\ell)_{\rm bottom}\right)
     =(-1)^{\rho(\rho-1)/2}Y_0^\rho.
\end{equation}
Only the reversing permutation can contribute: its column-index sum
already equals the minimum possible sum of nonzero entries in all
rows.  Thus the pencil has rank $\rho$.

At $s=\rho$, $Y_\ell$ is the formal variable $X_{\rho-\ell}$.
To obtain degree $\rho$ in the full determinant one must take all
these entries from the final block.  The complementary determinant
is \eqref{hier:Vcutoff}, which proves \eqref{hier:top}.
Its moments have degree at most $2L-2<h$, so
\eqref{hier:finite-moment} expresses each using finite rational
endpoints alone.  Positivity of the weight makes this $L\times L$
Gram matrix positive definite once at least $L$ distinct nodes have
been included.  This proves the last assertion, including the empty
case $L=0$.
\end{proof}

For example, $r=3,\rho=2,h=3$ gives the positive critical weight
\[
 R(x)=\frac16\left\{\frac1{(x+a)^2}-\frac3{(x+a+1)^2}
              +\frac3{(x+a+2)^2}-\frac1{(x+a+3)^2}\right\}.
\]
Its two coefficient matrices, indexed by harmonic orders, are
\[
 C^{(1)}=\begin{pmatrix}0&0&0\\0&0&0\\0&0&4\end{pmatrix},
 \qquad
 C^{(2)}=\begin{pmatrix}0&0&0\\0&0&1\\0&1&-4a-6\end{pmatrix}.
\]
The exact leading harmonic term is
$-\bigl(\sum_{n=0}^N R(n)\bigr)X_2^2$, and its coefficient is
nonzero for every $N\ge0$.

\subsection{Exact pole orders and the leading Laurent coefficients}

For an integer $m$ with $1\le m\le\rho$, put
\begin{equation}\label{hier:pole-parameters}
 \eta=\rho-m+1,\qquad L=r-\eta,
 \qquad \beta_m=\binom{h+m-1}{m-1}.
\end{equation}
Define the convergent lower-particle normalization
\begin{equation}\label{hier:lower-normalization}
 G_{L,h}^{(m)}(a)=
 \det\left(\sum_{n\ge0}n^{i+j}w_n^{[h]}(m,a)\right)_{0\le i,j<L},
 \qquad G_{0,h}^{(m)}=1.
\end{equation}

\begin{theorem}[A staircase of exact spectral poles]
\label{hier:poles}
For real $a>0$, $Z_{r,\rho}(s,a)$ has a pole of exact order
$\rho-m+1$ at every $m=1,\ldots,\rho$.  Its leading coefficient is
\begin{equation}\label{hier:leading-Laurent}
 \boxed{\displaystyle
 \lim_{s\to m}(s-m)^{\eta}Z_{r,\rho}(s,a)
 =(-1)^{\eta(\eta-1)/2}\beta_m^{\eta}G_{L,h}^{(m)}(a).}
\end{equation}
The lower normalization is strictly positive.  For complex $a$ with
$\Re a>0$, the same identity holds holomorphically, but at special
complex parameters the leading coefficient is not asserted nonzero.
\end{theorem}
\begin{proof}
Only $\zeta(s-m+1,a)$ has a pole at $s=m$; its residue is one.
Consequently the continued moment matrix has the local form
\[
 M(s)=\frac{C_{m-1}}{s-m}+A(s)
\]
with $A$ holomorphic.  The residue matrix is supported in its final
$\eta\times\eta$ block.  This block has zeros before the
anti-diagonal and constant anti-diagonal $\beta_m$, so its
determinant is
$(-1)^{\eta(\eta-1)/2}\beta_m^{\eta}$.

In the full determinant, the coefficient of $(s-m)^{-\eta}$ is this
residue determinant times the complementary top-left $L\times L$
minor of $A(m)$.  These top-left entries have no pole and coincide
with ordinary moments of $w_n^{[h]}(m,a)$.  Indeed, their largest
degree is $2L-2$ and the weight decays as $n^{-m-h}$; the required
strict inequality is
\[
 m+h-(2L-2)>1,
\]
which here has left side $\eta+2>1$.  The same inequality holds in
a small complex neighborhood, so normal convergence identifies the
analytic moment entries with their series there.  This proves
\eqref{hier:leading-Laurent}.  At positive real $a$, the Laplace
formula makes the weight strictly positive on infinitely many nodes,
so the lower Gram determinant is strictly positive.  The pole order
is therefore exact.
\end{proof}

\subsection{A Gamma product at the highest-order pole}

At $m=1$, put $L=r-\rho$.  The normalization in the preceding theorem
has the finite product
\begin{equation}\label{hier:Gamma-product}
 \boxed{\displaystyle
 G_{L,h}^{(1)}(a)=
 \prod_{j=0}^{L-1}
 \frac{(j!)^2\Gamma(a+j)\Gamma(h-L+1-j)}
      {\Gamma(h+1-j)\Gamma(a+h-j)}.}
\end{equation}
In particular, the leading Laurent coefficient at $s=1$ is
$(-1)^{\rho(\rho-1)/2}$ times this product.  For $\rho=r$ it is
the nonzero constant $(-1)^{r(r-1)/2}$.

Here is a derivation that also fixes the normalization.  The beta
partial-fraction identity gives
\[
 w_n^{[h]}(1,a)=\frac{\Gamma(n+a)}{\Gamma(n+a+h+1)}.
\]
Replace the two monomial bases of the $L$-particle moment determinant
by the monic bases $n^{\underline i}$ and $(n+a)_j$.  The mixed
moment is
\begin{equation}\label{hier:mixed-moment}
 M_{ij}=\frac{\Gamma(i+1)\Gamma(a+i+j)\Gamma(h-i-j)}
                  {\Gamma(h+1-j)\Gamma(a+h)},
 \qquad 0\le i,j<L.
\end{equation}
To see this, write $n=i+k$ and sum the resulting Gauss series with
parameters $i+1,a+i+j,a+i+h+1$.  Its excess is $h-i-j>0$;
in our parameters $h>2L-2$.  Thus the sum is justified by absolute
convergence.

After factoring rows and columns in \eqref{hier:mixed-moment}, the
remaining determinant is
\[
 \det_{0\le i,j<L}\frac{(a+i)_j}{(1-h+i)_j}.
\]
The elementary factorial determinant
\[
 \det\frac{(x_i)_j}{(x_i+t)_j}
 =\frac{\prod_{j=0}^{L-1}(t)_j\,
             \prod_{i<k}(x_k-x_i)}
            {\prod_{i=0}^{L-1}(x_i+t)_{L-1}}
\]
follows by multiplying each row by $(x_i+t)_{L-1}$ and comparing
polynomial coefficients; evaluation at $x_i=-i$ computes the
coefficient determinant.  Apply it with $x_i=a+i$ and $t=1-h-a$.
Use
\[
 (1-h-a)_j=(-1)^j\frac{\Gamma(a+h)}{\Gamma(a+h-j)},\qquad
 (1-h+i)_{L-1}=(-1)^{L-1}
                    \frac{\Gamma(h-i)}{\Gamma(h-i-L+1)}.
\]
All signs cancel, and the remaining factors give
\eqref{hier:Gamma-product}.  This is the same classical factorial
determinant mechanism as the predecessor's residue product; the new
point is its role as the coefficient of a pole of arbitrary order
$\rho$ in the shorter-stencil spectral family.

\subsection{A complete two-particle formula and its Stieltjes finite parts}

For $r=\rho=2$ one has $h=1$.  Directly from the three moment
identities,
\begin{equation}\label{hier:r2}
 \boxed{\displaystyle
 Z_{2,2}(s,a)=2a^{-s}\zeta(s-1,a)
        +(1-2a)a^{-s}\zeta(s,a)-\zeta(s,a)^2.}
\end{equation}
Thus, for every integer $m\ge3$, the ordinary convergent sum
\begin{align}\label{hier:r2-sum}
 &\sum_{0\le n<k}(k-n)^2
 \bigl((n+a)^{-m}-(n+a+1)^{-m}\bigr)
 \bigl((k+a)^{-m}-(k+a+1)^{-m}\bigr)\notag\\
 &\hspace{12mm}=2a^{-m}\zeta(m-1,a)
            +(1-2a)a^{-m}\zeta(m,a)-\zeta(m,a)^2.
\end{align}
For example, $a=1,m=3$ gives
$2\zeta(2)-\zeta(3)-\zeta(3)^2$.

Use the convention
\[
 \zeta(1+u,a)=\frac1u+
    \sum_{k\ge0}\frac{(-1)^k\gamma_k(a)}{k!}u^k,
 \qquad c=\frac1a-2.
\]
Expanding \eqref{hier:r2} at the double pole gives the full first
three Laurent terms
\begin{align}\label{hier:r2-s1}
 Z_{2,2}(1+u,a)
 ={}&-\frac1{u^2}+\frac{c-2\gamma_0(a)}u\notag\\
 &+2\gamma_1(a)-\gamma_0(a)^2
       +c\bigl(\gamma_0(a)-\log a+1\bigr)+O(u).
\end{align}
At the simple pole $s=2$,
\begin{align}\label{hier:r2-s2}
 \operatorname{FP}_{s=2}Z_{2,2}(s,a)
 ={}&2a^{-2}\bigl(\gamma_0(a)-\log a\bigr)
 +(1-2a)a^{-2}\zeta(2,a)-\zeta(2,a)^2.
\end{align}
These may be written entirely in generalized Stieltjes and polygamma
coordinates using $\gamma_0(a)=-\psi(a)$ and $\zeta(2,a)=\psi'(a)$.
They are spectral Laurent finite parts.  No equality with an
unsubtracted endpoint sum or a hard-cutoff constant is implicit.

There is a particularly simple check at order zero. Substitution of
$\zeta(0,a)=1/2-a$ and
$\zeta(-1,a)=-(a^2-a+1/6)/2$ into \eqref{hier:r2} gives
\begin{equation}\label{hier:zero-contact}
                  Z_{2,2}(0,a)=\frac1{12}
\end{equation}
for its spectral continuation. Yet at $s=0$ every weight
$(n+a)^0-(n+a+1)^0$ is exactly zero, so the unmodified defining
sum at that parameter has ordinary value zero. The discrepancy is
an explicit example of noncommutation between continuation and
specialization of the original infinite sum; \eqref{hier:zero-contact}
does not assign $1/12$ to an ordinary sum of zero terms.

All further Laurent coefficients of \eqref{hier:closed} are obtained
by finite determinant arithmetic from the Hurwitz expansions at the
relevant shifted orders.  In particular, at $s=m$ the coordinates are
the generalized Stieltjes constants from the sole resonant Hurwitz
factor, together with finitely many ordinary Hurwitz derivatives from
the nonresonant factors and elementary endpoint logarithms.  This is
an exact coefficient-generation procedure, without a claim that the
resulting higher constants admit an additional classical reduction.
